\section{Discussion and Outlook}
\label{sec:discussion_outlook}

Sections 2--5 establish the representative-level duality structure: a strong-weak involution on
the crossover coordinate and an orientation-gauge affine reflection for effective-temperature labels.
Section 6 demonstrates these mechanisms in the spin-boson qubit without elevating that model to a
claim of universality.

The transport mechanism requires two ingredients: (i) an explicit positive crossover coordinate with
distinct weak/strong asymptotics, and (ii) a closed representative map that acts algebraically on
the control labels. Whenever both are present, asymptotic work can be moved between branches without
changing the analytic functional family.

The proofs here are representative-level and asymptotic. They do not assert equality of physical
states under parameter inversion. They also do not provide a full classification of all models
admitting such discrete actions; the spin-boson section is an existence example with closed formulas.

Interpretation of the spin-boson example should remain anchored to established nonperturbative
results: Ohmic KT behavior and super-Ohmic delocalization are standard benchmarks, while sub-Ohmic
criticality is supported by NRG/QMC and complementary analytic approaches
\cite{leggettDynamicsDissipativeTwostate1987,weissQuantumDissipativeSystems2012,bullaQuantumPhaseTransitionSubohmic2003,vojtaQuantumPhaseTransitionsSubohmic2005,winterQuantumPhaseTransitionSubohmic2009,bullaNumericalRenormalizationGroup2008,chinGeneralizedPolaronAnsatz2011}.
In that context, the present contribution is a representative duality geometry for mean-force
calculations, not a replacement for full universality classification.

Natural next testbeds are multi-channel models, where multiple \(\chi_a\) coordinates may generate a
\((Z_2)^k\) reflection structure, and finite-\(N\) systems with non-trivial eigenlabel symmetries.
The minimal numerical protocol in Section 7 is designed to scale to those settings with only modest
changes.
